\section{결론}

본 연구는 웹 GUI pretraining에서 자연어 target을 어떻게 구성할지의
문제를 text-rich UI evidence에 대한 supervision coverage로
정식화한다. 동일한 화면과 구조화 evidence에서 Concise, Balanced,
Comprehensive verbalization을 만들고 Grounding-only 기준선과
비교한다. Stage 1 기반 coverage와 omission-targeted diagnostic,
공개 benchmark, PT부터 Trajectory까지의 checkpoint 평가를 연결해
어느 information budget이 웹 이해와 agent adaptation에 유용한지
검증한다. 이 설계는 긴 출력의 우월성을 전제하지 않으며,
포괄성의 이득과 target-token 비용ㆍ노이즈 사이의 균형을 밝히는
것을 목표로 한다.
